\chapter*{Resumen}
\addcontentsline{toc}{chapter}{Resumen}
\markboth{Resumen}{Resumen}

La enseñanza del Método de los Elementos Finitos (MEF) enfrenta una brecha pedagógica persistente: los textos clásicos exponen la formulación matricial con rigor pero en abstracto, y el software profesional opera como una caja negra que entrega la respuesta estructural sin el procedimiento que la produce, de modo que el estudiante aprende a operar sin formar criterio para interpretar y juzgar lo que obtiene. Esta tesis aborda esa brecha mediante el desarrollo de EduFEM, una herramienta educativa interactiva en español para el análisis estructural de problemas bidimensionales de elasticidad lineal en tensión plana y deformación plana, entendido como el análisis tenso-deformacional de medios continuos. La hipótesis sostiene que la exposición transparente e interactiva del canal de cálculo y de la respuesta estructural, sobre el modelo del propio estudiante, permitirá mejorar ese criterio. La herramienta combina un motor de cálculo con elementos isoparamétricos cuadriláteros Q4 y Q9, una interfaz gráfica que integra pre-proceso, proceso y post-proceso sobre un único lienzo de malla, ocho módulos educativos que exponen paso a paso cada etapa del cálculo, del mapeo isoparamétrico al ensamblaje global, y una memoria de cálculo automática con sustitución numérica que el estudiante puede reproducir y verificar a mano. La corrección del motor se sometió a una batería de verificación y validación: con el método de soluciones manufacturadas, en ambos estados planos y sobre mallas regulares y distorsionadas, se confirmaron las tasas de convergencia asintóticas teóricas del desplazamiento —orden dos en norma $L^2$ y orden uno en seminorma $H^1$ para Q4; tres y dos para Q9— y la convergencia del campo de tensiones recuperado; la viga de Timoshenko arrojó errores del 0,04\,\% en la tensión normal y del 0,26\,\% en la flecha frente a la solución analítica, y diferencias de hasta el 0,56\,\% frente a SAP2000; y la membrana de Cook evidenció el bloqueo por cortante (\emph{shear-locking}) del Q4 frente a la convergencia rápida del Q9, que a igual número de grados de libertad (GDL) reduce el error en un orden de magnitud. La validez de contenido del software se estableció con una tabla de especificaciones que contrasta sus instrumentos con las destrezas de interpretación y verificación de resultados que un consenso publicado de 67 expertos exige al egresado: las quince destrezas y los tres conceptos del universo tienen instrumento en EduFEM, y cada uno de los cuatro principios de aprendizaje que fundamentan el diseño tiene una encarnación comprobable. La hipótesis queda comprobada en su parte de diseño y contenido; su efecto sobre el criterio del estudiante se fundamenta en la literatura y se deja a una prueba de campo, diseñada en las recomendaciones.

\noindent\textbf{Palabras clave:} método de elementos finitos, software educativo, interpretación de resultados, elementos isoparamétricos Q4/Q9, verificación y validación, memoria de cálculo
